% Pre-ambulo
\documentclass[a4paper, 12pt]{abnt}

\usepackage[brazil]{babel}
\usepackage[T1]{fontenc}
\usepackage[utf8]{inputenc}
\usepackage{dsfont}
\usepackage{amssymb,amsmath}
\usepackage{multirow}
\usepackage[alf]{abntcite}
\usepackage[pdftex]{color, graphicx}
\usepackage{colortbl}
\usepackage{url}
\usepackage{abnt-alf}
\usepackage{abntcite}
\usepackage{algorithm}
\usepackage{algorithmic}

%\usepackage{alg}


% Redefinicao de instrucoes
\floatname{algorithm}{Algoritmo}
\renewcommand{\algorithmicrequire}{\textbf{Entrada:}}
\renewcommand{\algorithmicensure}{\textbf{Saída:}}
\renewcommand{\algorithmicend}{\textbf{fim}}
\renewcommand{\algorithmicif}{\textbf{se}}
\renewcommand{\algorithmicthen}{\textbf{então}}
\renewcommand{\algorithmicelse}{\textbf{senão}}
\renewcommand{\algorithmicfor}{\textbf{para}}
\renewcommand{\algorithmicforall}{\textbf{para todo}}
\renewcommand{\algorithmicdo}{\textbf{faça}}
\renewcommand{\algorithmicwhile}{\textbf{enquanto}}
\renewcommand{\algorithmicloop}{\textbf{loop}}
\renewcommand{\algorithmicrepeat}{\textbf{repetir}}
\renewcommand{\algorithmicuntil}{\textbf{até que}}
\renewcommand{\algorithmiccomment}[1]{\% #1}


\newcommand{\ignore}[1]{}


% Hifenização de palavras feita de forma incorreta pelo LaTeX
\hyphenation{PYTHON ou-tros}


% Inicio do documento
\begin{document}

	\frenchspacing

	% Capa (arquivo Includes/Capa.tex)
	% Capa
% Proteção externa do trabalho e sobre a qual se imprimem as informações indispensáveis
% à sua identificação.

% Especificação da capa
\begin{titlepage}
	\begin{center}

		% Cabeçalho (não deve ser modificado)
		% Contém o brasão da Universidade, o logotipo do Departamento, além dos dados
		% relacionados à vinculação do aluno (Universidade, Centro, Departamento e Curso)
		\begin{minipage}{2cm}
			\begin{center}
				\includegraphics[width=1.7cm, height=2.0cm]{Imagens/Brasao-UFRN.jpg}
			\end{center}
		\end{minipage}
		\begin{minipage}{11cm}
			\begin{center}
				\begin{espacosimples}
					{\small \textsc{Universidade Federal do Rio Grande do Norte}			\\
							  \textsc{Centro de Ciências Exatas e da Terra}						\\
							  \textsc{Departamento de Informática e Matemática Aplicada}	\\
							  \textsc{Bacharelado em Engenharia de Software}}
				\end{espacosimples}
			\end{center}
		\end{minipage}
		\begin{minipage}{2cm}
			\begin{center}
				\includegraphics[width=1.8cm, height=1.5cm]{Imagens/Logotipo-DIMAp.jpg}
			\end{center}
		\end{minipage}

		\vspace{5cm}

		% Título do trabalho
		{\setlength{\baselineskip}%
		{1.3\baselineskip}
		{\LARGE \textbf{Twydi: Uma Ferramenta de Documentação para Equipes de Desenvolvimento de Software}}\par}

		\vspace{4cm}

		% Nome do aluno (autor)
		{\large \textbf{Igor Marques da Silva}}

		\vspace{7cm}

		% Local da instituição onde o trabalho deve ser apresentado e ano de entrega do mesmo
		Natal-RN\\Novembro 2015
	\end{center}
\end{titlepage}


	% Folha de rosto (arquivo Includes/FolhaRosto.tex)
	% Folha de rosto
% Contém os elementos essenciais à identificação do trabalho.

% Título, nome do aluno e respectivo orientador e filiação
\titulo{\Large{Estudo e Implementação de uma Ferramenta de Documentação em Suporte à Reuso de Conhecimento em Equipes de Desenvolvimento de Software}}
\autor{Igor Marques da Silva}
\orientador[Orientador]{\par Prof. Dr. Fernando Marques Figueira Filho}
\instituicao
{
	Universidade Federal do Rio Grande do Norte -- UFRN \par
	Departamento de Informática e Matemática Aplicada -- DIMAp
}

% Natureza do trabalho (não deve ser modificada)
\comentario
{
	Proposta de Monografia de Graduação apresentada ao Departamento de Informática e Matemática Aplicada do
	Centro de Ciências Exatas e da Terra da Universidade Federal do Rio Grande do Norte como
	requisito parcial para a obtenção do grau de bacharel em Engenharia de Software.
}

% Local e data
\local{Natal-RN}
\data{Novembro de 2015}

\folhaderosto

	% Folha de aprovacao (arquivo Includes/FolhaAprovacao.tex)
	% Folha de aprova��o
\begin{folhadeaprovacao}
	\setlength{\ABNTsignthickness}{0.4pt}
	\setlength{\ABNTsignwidth}{10cm}
	
	% Informa��es gerais acerca do trabalho 
	% (nome do autor, t�tulo, institui��o � qual � submetido e natureza)
	\noindent 
	Monografia de Gradua��o sob o t�tulo \textit{T�tulo da monografia} apresentada por 
	Nome do aluno e aceita pelo Departamento de Inform�tica e Matem�tica Aplicada do
	Centro de Ci�ncias Exatas e da Terra da Universidade Federal do Rio Grande do Norte,
	sendo aprovada por todos os membros da banca examinadora abaixo especificada:
		
	% Membros da banca examinadora e respectivas filia��es
	\assinatura
	{
		Titula��o e nome do(a) orientador(a)\\
		{\small Orientador(a)} 															\\ 
		{\footnotesize
			Departamento 																	\\
		  	Universidade
		}
	}
	
	\assinatura
	{
		Titula��o e nome do membro da banca examinadora							\\
		{\small Co-orientador(a), se houver}										\\ 
		{\footnotesize
			Departamento 																	\\
		  	Universidade
		}
	}
		
	\assinatura
	{
		Titula��o e nome do membro da banca examinadora 						 \\ 
		{\footnotesize
			Departamento 																	 \\
		  	Universidade
		}
	}
		
	\assinatura
	{
		Titula��o e nome do membro da banca examinadora 						 \\ 
		{\footnotesize
			Departamento 																	 \\
		  	Universidade
		}
	}
		
	\vfill
	
	\begin{center}
		Natal-RN, data de aprova��o (por extenso).
	\end{center}
\end{folhadeaprovacao}

	% Dedicatoria (arquivo Includes/Dedicatoria.tex)
	% Dedicat�ria

\chapter*{}
\vspace{15cm}
\begin{flushright}
	Homenagem que o autor presta a uma ou mais pessoas.
\end{flushright}

	% Agradecimentos (arquivo Includes/Agradecimentos.tex)
	% Agradecimentos

\chapter*{Agradecimentos}

Em quatro anos de graduação, muito amadureci e muito aprendi. Sem dúvida tudo isso não seria possível sem uma boa parcela de companheirismo, apoio e dedicação. Tudo o que alcancei foi graças a uma boa parcela de pessoas.

Obrigado à minha família, por tanto apoio. Não seria nada sem vocês. Literalmente. Obrigado por dia após dia me aturarem, me educarem e me proverem com tudo que há de melhor nesse mundo.

Obrigado à minha companheira de todas as horas. Obrigado por tanta paz, tanto carinho e tanto amor todos os dias.

Agradeço também a todos os meus amigos de ensino médio da querida turma 401 do IFRN, porque vocês merecem mesmo não tendo nada a ver com esse trabalho.

Devo agradecer também a todos que me ajudaram a adquirir tantos ensinamentos ao longo desses quatro anos. Sei o que sei hoje da área de computação graças a uma grande parcela de amigos, colegas e professores.

Obrigado especialmente à 4Soft e a todos que fizeram ou fazem parte dela. Não seria o profissional que sou hoje sem essa tão caristmática (e problemática) empresa júnior. Foram três anos e meio de muita luta, aprendizados e ódio acumulado (por fatores que não vem ao caso) e que sem dúvida abriram meus olhos com relação ao mundo.

Um agradecimento especial também aos meus companheiros de curso que mais me ensinaram nesse tempo todo (e que certamente continuarão ensinando). Vou citá-los em ordem alfabética para que ninguém fique com ciúmes, ok? Obrigado Bernardo, Iago, Lucas, Luiz Rogério e Waldyr por tantas dicas trocadas, \textit{pulls requests} revisados, dúvidas esclarecidas e trabalhos feitos.

Agradeço também a todos os professores do Departamento de Informática e Matemática Aplicada da UFRN, em especial ao professor Fernando, pelas infindáveis orientações e revisões deste trabalho. Obrigado também ao professor Nélio Cacho por ter passado o trabalho mais desafiador na disciplina mais insana que paguei durante a graduação e que serviu de um amadurecimento enorme.

Obrigado também a qualquer um que se sinta válido de estar mencionado aqui mas acabou não sendo por algum motivo qualquer. Desculpa.

Obrigado a todos que torceram por mim.

Obrigado.


   % Epigrafe (arquivo Includes/Epigrafe.tex)
	% Ep�grafe (cita��o seguida de indica��o de autoria)

\chapter*{}
\vspace{15cm}
\begin{flushright}
	\textit
	{
		Cita��o
	}\medskip\\ 
	Autor
\end{flushright}

	% Resumo em língua vernacula (arquivo Includes/Resumo.tex)
	% Resumo em língua vernácula
\begin{center}
	{\Large{\textbf{Twydi: Uma Ferramenta de Documentação para Equipes de Desenvolvimento de Software}}}
\end{center}

\vspace{1cm}

\begin{flushright}
	Autor: Igor Marques da Silva\\
	Orientador: Prof. Dr. Fernando Marques Figueira Filho
\end{flushright}

\vspace{1cm}

\begin{center}
	\Large{\textsc{\textbf{Resumo}}}
\end{center}

\noindent
Desenvolvedores de software constantemente trocam informações entre si para o aprimoramento contínuo de suas habilidades e resolução de problemas. Muitas dessas informações são perdidas devido ao uso de mídias que não oferecem bom suporte ao armazenamento, recuperação e visualização de informações.

Este trabalho apresenta uma ferramenta que permite aos desenvolvedores de agregar descrições textuais, código, referências externas em artefatos, além de possibilitar a recuperação destes, bem como os estudos realizados ao longo do seu processo de desenvolvimento e avaliação. Apesar de melhorias no quesito de pré-visualização dos artefatos serem necessárias, a ferramenta atende aos requisitos de suporte a troca de conhecimento de desenvolvedores através do armazenamento do registro de soluções e compartilhamento de informações e referências para a equipe. Os estudos realizados neste trabalho, bem como as lições aprendidas durante o desenvolvimento da ferramenta desenvolvida em si, servem como base para trabalhos futuros de melhoria da qualidade de documentação e processos de troca de conhecimento em equipes de desenvolvimento de software.

\noindent\textit{Palavras-chave}: Documentação de Software. Reuso de Software. Desenvolvimento de Software.


	% Abstract, resumo em língua estrangeira (arquivo Include/Abstract.tex)
	% Resumo em l�ngua estrangeira (em ingl�s Abstract, em espanhol Resumen, em franc�s R�sum�)
\begin{center}
	{\Large{\textbf{T�tulo do trabalho (em l�ngua estrangeira)}}}
\end{center}

\vspace{1cm}

\begin{flushright}
	Author: Nome do aluno\\
	Advisor: Titula��o e nome do(a) orientador(a)
\end{flushright}

\vspace{1cm}

\begin{center}
	\Large{\textsc{\textbf{Abstract}}}
\end{center}

\noindent O resumo em l�ngua estrangeira (em ingl�s \textit{Abstract}, em espanhol \textit{Resumen}, em franc�s \textit{R�sum�}) � uma vers�o do resumo escrito na l�ngua vern�cula para idioma de divulga��o internacional. Ele deve apresentar as mesmas caracter�sticas do anterior (incluindo as mesmas palavras, isto �, seu conte�do n�o deve diferir do resumo anterior), bem como ser seguido das palavras representativas do conte�do do trabalho, isto �, palavras-chave e/ou descritores, na l�ngua estrangeira. Embora a especifica��o abaixo considere o ingl�s como l�ngua estrangeira (o mais comum), n�o fica impedido a ado��o de outras linguas (a exemplo de espanhol ou franc�s) para reda��o do resumo em l�ngua estrangeira.

\noindent\textit{Keywords}: Keyword 1, Keyword 2, Keyword 3.

	% Lista de figuras
% 	\listoffigures

	% Lista de tabelas
% 	\listoftables

	% Lista de abreviaturas e siglas
% 	\listadeabreviaturas

	% Lista de símbolos
% 	\listadesimboenlos

	% Lista de algoritmos (se houver)
	% Devem ser incluídos os pacotes algorithm e algorithmic
	% \listofalgorithms

	% Sumário
	\sumario

	% Parte central do trabalho, englobando os capítulos que constituem o mesmo
	% Os referidos capítulos devem ser organizados dentro do diretório "Capítulos"

	\chapter{Introdução}

Desenvolvedores de software são trabalhadores do conhecimento e frequentemente precisam reutilizar o conhecimento aprendido por outros desenvolvedores, de modo a proporcionar um melhor desempenho para toda organização~\cite{Druker1993}~\cite{Wiig2003}.
É esperado que os trabalhadores do conhecimento melhorem  continuamente seu trabalho em um processo em culmina na melhora significativa da sua empresa~\cite{Kavitha2011}.

% IM: Fernando, posso passar esses dois parágrafos abaixo pra revisão?

Em equipes de desenvolvimento de software, o reuso de conhecimento ocorre utilizando uma variedade de artefatos, como código fonte, requisitos, modelos, dados e padrões~\cite{Levy2009}, bem como através de interações face a face, comunicação escrita e também via o repasse de referências de documentação, links, dentre outros~\cite{Storey2014}~\cite{Olson2000}~\cite{CubraniC2004}.

Neste quesito, entra a gerência de conhecimento, encarregada da elicitação, armazenamento, gerenciamento e reuso do conhecimento em suas diferentes formas.
\ignore{FONTE 2}.
Dentre esse reuso, se encontra, por exemplo, como se deu a implementação de uma determinada funcionalidade em um determinado projeto de software.

% FONTE 2 - KM is comprised of the elicitation, packaging and management, and reuse of knowledge in all of its different forms, and in particular, software engineering artifacts as code, design, requirements, models, data, standards, and lessons learned.
%Software development can be improved by recognizing the related knowledge content and structure as well as the required knowledge, and performing planning activities. - 05071412.pdf

% IM: Fernando, não consigo achar esse artigo da FONTE 2, vc poderia me passar por favor?
% IM: Seria esse? http://www.computer.org/csdl/proceedings/chase/2009/3712/00/05071412-abs.html

A troca de informações é um mecanismo fundamental para que o reuso de conhecimento ocorra de maneira eficiente em equipes de desenvolvimento. Desenvolvedores em fase de aprendizado de uma determinada tecnologia comumente passam por situações e problemas semelhantes. Quando se deparam com algum tipo de adversidade, procuram por alguma fonte de informação capaz de auxiliá-los por tal problema. Existe então o desperdício de tempo de se recuperar tal informação (muitas vezes, interferindo nas atividades de um colega de trabalho) para se resolver um problema que já é de conhecimento da equipe dado alguma experiência anterior.

% IM: Fernando, posso botar o parágrafo abaixo para a revisão?

Com a falta de uma boa gerência de conhecimento, é comum desenvolvedores implementarem funcionalidades semelhantes em diferentes contextos usando abordagens ad-hoc~\cite{SangMok2011}. Essas novas implementações, principalmente se feitas por desenvolvedores menos experientes, tendem a ser menos otimizadas e elegantes que soluções já previamente implementadas pela equipe, podendo gerar um débito técnico a longo prazo.

Desenvolvedores mais experientes em um determinado projeto tendem a atuar como mentores~\cite{CubraniC2004} e tal ato, como atividade de gerência de conhecimento, tende a consumir recursos~\cite{Wiig2003}. Um dos pilares da gerência de conhecimento é a melhoria da produtividade através do compartilhamento e transferência eficiente de conhecimento, que tende a consumir tempo e as vezes até impossível de se realizar~\cite{Levy2009}.
Muitas das interações entre desenvolvedores são informais e não há registro que possa ser consultado para propocionar o reuso do conhecimento trocado~\cite{Olson2000}. Um agravante tipicamente é a frequente rotatividade de membros em equipes. Nesses casos, a saída de um membro da equipe que detém determinado conhecimento pode prejudicar toda organização. De fato, a falta de uma abordagem mais sistemática para proporcionar o reuso de conhecimento pode estar associada ao fracasso de projetos em organizações~\cite{Hall2008}.

Assim, a elaboração de uma ferramenta capaz de agregar referências de código e tarefas a soluções pode trazer enormes benefícios a equipes de desenvolvimento~\cite{CubraniC2004}. Este estudo propôs a elaboração de tal ferramenta. Esta atua como um catálogo, agregando referências de código, informações externas e comentários fornecidos por desenvolvedores da equipe com o intuito de auxiliar outros desenvolvedores a buscar em fontes da própria equipe como se deram implementações de funcionalidades já feitas em outros projetos.

Segue um exemplo de uso da ferramenta: Um determinado desenvolvedor X, experiente dentro de sua equipe, percebe que muitos membros novatos procuram sua ajuda para implementar uma funcionalidade de exportação de uma página web para um documento em formato PDF. X então pode fazer uso da ferramenta, para registrar como se realiza a implementação de tal funcionalidade solicitada. Ele informa um título, descrição curta (para facilitar futuras buscas) e uma descrição de como se dá tal implementação. Como ilustração em formato de código, o desenvolvedor pode utilizar de links de trechos de código disponívels artefatos de código já implementados pela sua equipe. Esses links são renderizados no editor de texto, sem a necessidade de X copiar e colar o código em si dentro do editor. X também pode complementar sua descrição adicionando links para outras páginas web (respostas de sites como Stack Overflow, outros tutoriais, etc...) em forma de anexo da documentação. Por fim, X informa tags relevantes ao artefato que está produzindo, facilitando a sua recuperação por parte de outros membros de sua equipe.

Outros desenvolvedores, a partir de então, podem recorrer diretamente a ferramenta quando precisarem implementar a funcionalidade de exportação de página em formato PDF. Dessa forma, X poderá se dedicar mais a outras atividades de seu dia-a-dia de trabalho.

Este estudo descreve a elaboração de tal ferramenta e sua avaliação em dois contextos. O primeiro é o de uma equipe real de desenvolvimento de software e o segundo contexto é o de alunos de graduação da área de computação enquanto cursavam a disciplina de Desenvolvimento Colaborativo de Software. A ferramenta foi utilizada para se auxiliarem durante o desenvolvimento do projeto final de disciplina.

O objetivo geral deste trabalho é responder as seguintes perguntas:

\begin{enumerate}
  \item A ferramenta diminui o esforço gasto com mentoria de desenvolvedores?
  \item Desenvolvedores preferem utilizar a ferramenta aos métodos tradicionais de aprendizado/ resolução de dúvidas?
  \item Quais atributos da ferramenta oferecem vantagens da mesma em relação aos métodos tradicionais de aprendizado/ resolução de dúvidas?
\end{enumerate}

% IM: To falando muito "A ferramenta"...

% O estudo também realiza uma análise crítica dos impactos do uso da mesma no contexto de empresa. Dentre os impactos previstos estão a redução da necessidade de mentoria ou consulta para os casos de  reuso de funcionalidades em diferentes contextos. Tais impactos visam contribuir para a redução de gastos de recursos relacionados a tempo de implementação de funcionalidades, esforço de recuperação de informações, orientação de colaboradores e trabalho dos desenvolvedores que atuam como mentores dentro da equipe.
% IM: Vou tirar esse parágrafo (acima) por ter botado as perguntas de pesquisa

Para responder as perguntas de pesquisa propostas, foi feito aplicado um questionário com usuários de ferramenta, além de entrevistas semi-estruturadas com alguns dos usuários.

Concluiu-se que...
% IM: Vou precisar dizer os resultados aqui, né, professor?

O Capítulo 2 cita as referências utilizadas neste estudo como o uso de mídias pelo engenheiro de software, os efeitos da rotatividade em equipes de software, além de gestão e recuperação de informação. O Capítulo 3 explica em detalhes os procedimentos tomados para a construção e análise da ferramenta. O Capítulo 6 descreve em detalhes a ferramenta elaborada. O Capítulo 5 realiza a análise dos resultados dos experimentos bem como discute seus resultados. O Capítulo 6 aborda as considerações finais sobre o trabalho.
% IM: tá muito "crua"?

% FF: algumas coisas para você incluir aqui:
% 1- como será feita esta análise? que métodos de avaliação serão utilizados? (colocar de forma bem resumida)
% IM: Gostaria de conversar sobre este ponto com vc depois.

% 3- de onde veio o texto abaixo em inglês? ele está melhor que o em português acima.
% IM: Acredito que o texto veio de algum slide que fizemos. Mas coloquei uma descrição de um caso de uso acima. Fiz também uma descrição geral da ferramenta na linha 37. Está vaga demais?

% In software development, it is common for a member of the team to implement a feature very similar to one that was already done by a teammate. In some cases, that is a cumbersome task. Searching through documentation and code, as well as asking other teammates can be time and resource demanding.
% We propose a tool that will permit a team to store documentation for features via GitHub issues and comm its merged with other information like Stack Overflow questions, framework or language documentation or even notes written by teammates. This tool will be designed, tested and evaluated by the developers of 4Soft, the Junior Enterprise of Software Engineer of UFRN.

% FF: coloque um .gitignore pra LaTeX e arquivos relacionados? :)
% IM: Não entendi :/. Se eu colocar um gitignore pra arquivos relacionados não vamos ter como compartilhar um com o outro via git. Além disso, o overleaf não permite o upload de PDFs dentro da estrutura dos documentos, por exemplo. Posso compartilhar uma pasta de referencias minhas no dropbox com vc, que tal?


	\chapter{Objetivos}

% gerais e específicos
Este trabalho tem por finalidade analisar que influência uma ferramenta de documentação de funcionalidades de software traz para equipes de desenvolvimento. Além disso, objetiva verificar como tal ferramenta pode contribuir para a redução significativa de tempo despendido para a explicação repetitiva de tarefas já executadas anteriormente e como a troca de conhecimento entre membros pode fluir de maneira melhor e com baixo custo.

% A ferramenta está prevista de possuir os seguintes requisitos:

% \begin{enumerate}
% \item Gerenciamento de catálogo de funcionalidades implementadas em um projeto com integração com repositório de código no GitHub
% \item Anexo de fontes externas que auxiliem a implementação de determinada funcionalidade
% \item Anexo de descrição de etapas da implementação de determinada funcionalidade
% \item Recuperação de funcionalidades no catálogo
% \item Anexo de variações da implementação de determinada funcionalidade
% \end{enumerate}

% Assim, tal ferramenta prevê inicialmente os benefícios de fácil recuperação de referências sem a necessidade de intermédio de colegas de trabalho e, por consequência, garantindo que menos tempo será demandado para a orientação de desenvolvedores.

Este trabalho visa então responder às seguintes perguntas de pesquisa:

\begin{enumerate}
\item Quais são os requisitos para o suporte ferramental em apoio à transferência e reuso do conhecimento organizacional em empresas de desenvolvimento de software?
\item Que outras ferramentas existentes oferecem tal suporte?
\item Como estimular a adoção e o uso desse suporte ferramental em equipes de desenvolvimento de software? 
\item Que benefícios e limitações tal suporte provê a equipes de desenvolvimento?

\end{enumerate}



	\chapter{Revisão de Literatura}

\section{Uso das mídias pelo engenheiro de software}

A atividade de desenvolvimento de software requer que seus profissionais se mantenham constatemente atualizados, devido ao fluxo constante de inovações que surgem diariamente ~\cite{Singer2014}.

Uma das formas que desenvolvedores utilizam tal são as redes sociais ~\cite{Treude2012}~\cite{Storey2014} além de seus próprios colegas de trabalho, tendo em vista que desenvolvedores frequentemente recorrem a própria equipe para obter ajuda~\cite{Weinberg1998}.

% Fred Brooks: The Mythical Man-Month
% IM: não consegui achar algo sobre aprendizado dentro da própria equipe nele.

Sobre as mídias em geral, Clark e Brennan (1991) estabeleceram que estas podem apresentar as seguintes características:

\begin{itemize}
  \item Copresença: está no mesmo ambiente físico
  \item Visibilidade: está visível aos participantes da comunicação
  \item Contemporaneidade: a mensagem é recebida imediatamente
  \item Simultaneidade: Ambos os participantes podem enviar e receber mensagens
  \item Sequencialidade: turnos não podem sair de sequencia
  \item Revisibilidade: permite rever as mensagens enviadas
  \item Revisabilidade: permite rever uma mensagem antes de ser enviada
\end{itemize}

A comunicação com colegas de trabalho face-a-face comparado a outras mídias, provê uma série de vantagens, conforme apresentado por Olson et al (2000)(?) no quadro abaixo.

https://www.dropbox.com/s/3aq85imoew7vidv/Screenshot%202015-10-27%2013.59.20.png?dl=0

Porém, não apresenta uma importante características para a gestão do conhecimento de uma equipe: a revisibilidade.

Ferramentas de \textit{chat}, por sua vez, apresentam esta caracterísica, mas ainda sim possuem problemas de recuperação fácil de informações trocadas nela (variando de ferramenta para ferramenta).

\section{Rotatividade de pessoal}

Rotatividade de pessoal é inerente a qualquer área de atuação da indústria [CARECE DE FONTES].
Sua ocorrência acarreta em [PROBLEMAS].
% staff turnover (software development), Employee turnover, turnover open source development

Em equipes de desenvolvedores de software, a partida de um membro pode acarretar em [FALAR DE COISAS COMO PERDA DE INFORMAÇÕES] caso não haja nenhuma forma de registro em mídia da mesma.
% knowledge management, tacit knowledge, software development (organization or team)

Em geral, novos membros da equipe passam por um período de adaptação onde devem aprender padrões, práticas e ferramentas utilizadas pelo grupo [CARECE DE FONTES].
% 5 - PEDIR A FERNANDO - Rabelo2015
% Knowledge management and organizational culture in a software organization: a case study
% https://scholar.google.com/citations?view_op=view_citation&hl=en&user=I8o8rfoAAAAJ&sortby=pubdate&citation_for_view=I8o8rfoAAAAJ:maZDTaKrznsC


% 6 - Social Barriers Faced by Newcomers - Steinmacher2015

Em casos específicos, como desenvolvedores iniciantes (em geral chamados de \"desenvolvedores júnior\", no mercado), ainda existe a necessidade em se adquirir conhecimento sobre as tecnologias utlizadas pela equipe ao qual foi recém integrado.

Todo esse conhecimento é em geral transmitido oralmente [CARECE DE FONTES] ou por mídias que não oferecem um bom suporte à sua manutenção [CARECE DE FONTES].
% neste parágrafo você explicita quais são os problemas das abordagens atuais. ao mesmo tempo, você justifica o seu trabalho.

Sendo assim, uma ferramenta que estimule seus usuários a registrarem informações fundamentais para a equipe tende a evitar tanto que informações sejam perdidas no caso de partidas quanto facilitar o aprendizado e adaptação de novos membros à equipe.

\section{Gestão e recuperação de informação}

A gestão de informações é a área [ARRUMAR DEFINIÇÃO].
% 5 - PEDIR A FERNANDO - Rabelo2015
% Knowledge management and organizational culture in a software organization: a case study


Uma boa gestão de informação e facilidade de se recuperar informações permitem que [ARRUMAR BENEFÍCIOS].

É possível ver então que uma equipe de desenvolvimento que apresente alguma forma de gestão de informações pode garantir facilitar e estimular, por exemplo, o reuso de soluções já conhecidas de problemas recorrentes.

% \section{GitHub}

% 7 - A study on the geographical distribution - Figueira2015

\section{Reuso em engenharia de software}

Reuso é uma das caracterísitcas chaves de um código de qualidade. Sua definição é [ARRUMAR DEFINIÇÃO DE REUSO].
% ver o livro da Gang of Four

Bons engenheiros de software procuram desenvolver código com boa possibilidade de reuso, caso seja conveniente [CARECE DE FONTES].

É notável então que uma boa forma de garantir a gerência de informações visando o reuso de conteúdo produzido por uma equipe pode trazer melhorias para a qualidade do código produzido pela mesma.
   

	\chapter{Métodos}

% O que será feito, como será feito, … para abordar o problema

Este trabalho prevê o desenvolvimento da ferramenta de documentação de implementação funcionalidades. Nela, será possível vincular recursos do repositório do projeto no GitHub\footnote{\url{www.github.com}} (\textit{issues}, \textit{commits}, etc) a requisitos e suas implementações, bem como outros referenciais (\textit{links} para perguntas no Stack Overflow\footnote{\url{www.stackoverflow.com}}, desenhos, por exemplo) de forma a gerar um guia ou tutorial de como realizar tal implementação novamente no futuro.

A empresa 4Soft terá participação significativa em todo o estudo, desde a concepção até no uso da ferramenta.

% metodo: entrevista, questionário de satisfacao. Focar em pensar em como responder as perguntas

As etapas do estudo são:


\begin{description}
   \item Estudo de aplicações existentes

		Será feita uma busca por aplicações que realizem atividades semelhantes às propostas. Suas limitações serão analisadas pela equipe de pesquisa e um panorama inicial será traçado de modo que a ferramenta proposta possa suprir as necessidades iniciais e as limitações encontradas.
    
    \item Inquérito contextual

		Entrevistas e seções de \textit{brainstorm} serão feitas com os participantes da empresa júnior mencionada. Será analisado como se dá seu processo de trabalho e como pode se dar o fluxo de atividades na ferramenta. As informações necessárias serão coletadas através de entrevistas e aplicação de questionários de satisfação.

	\item Sessões de interpretação da equipe

		Reuniões com a equipe de pesquisadores que trabalharão no projeto serão feitas para definir com mais detalhes o escopo da ferramenta, bem como será sua arquitetura e implementação.

	\item Prototipação e implementação da ferramenta

		Nesta etapa, inicialmente, protótipos de baixa fidelidade serão elaborados pelo pesquisador. Posteriormente, serão expostos a todos os participantes do projeto (empresa júnior e pesquisadores) e seu \textit{feedback} será colhido e analisado. A partir daí, a aplicação passará para a etapa de implementação seguindo processo iterativo e incremental de desenvolvimento de software.

	\item Implantação e observação do uso da ferramenta

		A ferramenta então estará disponível para uso de todos os membros da empresa 4Soft. A adesão dos desenvolvedores a ferramenta será analisada nesta etapa, bem como seu uso monitorado (quantidade de artefatos de documentação criados, por exemplo).
        
        
	\item Avaliação da ferramenta

		Ao fim do período anterior, uma nova bateria de entrevistas e aplicação de questionários de satisfação serão realizadas para avaliar qualitativamente como se deu a utilização da ferramenta, bem como se deram os efeitos de seu uso e a satisfação de seus usuários. 
		A versão final do trabalho com os resultados será então redigida.


\end{description}

	\chapter{Plano de Trabalho}

% Tarefas ou Atividades (bem macro)
% Cronograma

O cronograma de atividades se dará conforme o seguinte planejamento:

\begin{enumerate}

\item Estudo de aplicações existentes: Dezembro/2014

\item Inquérito contextual: Janeiro/2015

\item Sessões de interpretação da equipe: Janeiro/2015

\item Prototipação e implementação da ferramenta: Janeiro/2015 a Maio/2015

\item Implantação e observação do uso da ferramenta: Junho/2015 a Setembro/2015

\item Avaliação da ferramenta: Outubro/2015 a Novembro/2015

\end{enumerate}



	% Bibliografia (arquivo Capitulos/Referencias.bib)
	\bibliography{Capitulos/Referencias}
	\bibliographystyle{abnt-alf}

	% Página em branco
	\newpage

\end{document}